\section{The proper-smooth parity question and the rational spectral
target}\label{the-proper-smooth-parity-question-and-the-rational-spectral-target}

\subsection{Exact known results answering the parity
question}\label{exact-known-results-answering-the-parity-question}

Deligne writes, immediately after Corollary 4.1.5:

\begin{quote}
Ce résultat devrait valoir sans hypothèse de projectivité, mais je ne
sais pas le prouver en caractéristique \(\ne0\).
\end{quote}

The source is \emph{Weil II}, printed p.~218.

The question following Deligne's Corollary 4.1.5 in \emph{Weil II},
printed p.~218, concerns removing projectivity from the odd Betti-number
parity statement in the \textbf{proper smooth} setting. That
interpretation is explicitly identified by Sun and Zheng (2016, p.~236).
The statement is now known. Two results must be distinguished:

\begin{longtable}[]{@{}
  >{\raggedright\arraybackslash}p{(\columnwidth - 4\tabcolsep) * \real{0.3333}}
  >{\raggedright\arraybackslash}p{(\columnwidth - 4\tabcolsep) * \real{0.3333}}
  >{\raggedright\arraybackslash}p{(\columnwidth - 4\tabcolsep) * \real{0.3333}}@{}}
\toprule\noalign{}
\begin{minipage}[b]{\linewidth}\raggedright
Statement
\end{minipage} & \begin{minipage}[b]{\linewidth}\raggedright
Exact primary result
\end{minipage} & \begin{minipage}[b]{\linewidth}\raggedright
Scope
\end{minipage} \\
\midrule\noalign{}
\endhead
\bottomrule\noalign{}
\endlastfoot
Odd-degree Betti numbers are even & Suh (2012), Corollary 2.2.3, p.~298
& Any proper smooth variety over any field; cohomology over a separably
closed extension; every odd degree \(r\ge1\) \\
Each of \(+q^{r/2}\) and \(-q^{r/2}\) has even algebraic multiplicity on
\(H^r\) for odd \(r\) & Suh (2012), Theorem 3.3.1, p.~301, applied using
proper-smooth purity; independently and more strongly, Sun--Zheng
(2016), Corollaries 1.2 and 1.4, pp.~236--237 & Proper smooth
\(X/\mathbf F_q\), without a projectivity or semisimplicity
assumption \\
\end{longtable}

Suh's Theorem 3.3.1 is stated for every separated scheme of finite type
over \(\mathbf F_q\): each of the two real numbers \(\pm q^{r/2}\),
\(r\ge1\) odd, has even multiplicity as a reciprocal zero or pole of the
zeta function. For a proper smooth scheme, purity of \(H^j\) isolates
weight \(r\) in degree \(r\); thus there is no cancellation with another
degree, and the theorem gives the second row. For a general scheme, its
zeta-function statement does not by itself give the same assertion in
each individual cohomological degree.

The later theorem supplies an actual equivariant form. Sun--Zheng,
Theorem 1.1, p.~236, states that for a proper equidimensional scheme
over an arbitrary field, its \(n\)-th intersection cohomology has a
Galois-equivariant perfect pairing with values in
\(\mathbf Q_\ell(-n)\), symmetric for even \(n\) and alternating for odd
\(n\). Their \textbf{Corollary 1.2 on the same page} gives the direct
ordinary-cohomology statement for every proper smooth scheme: \[
 H^n(X_{\bar k},\mathbf Q_\ell)\otimes
 H^n(X_{\bar k},\mathbf Q_\ell)
       \longrightarrow\mathbf Q_\ell(-n).
                                                        \tag{CC-XC.P33}
\] This is stronger than evenness of the total odd Betti number.

Over \(\mathbf F_q\), let \(\mu_{\lambda,e}\) count the Jordan blocks of
size \(e\) with eigenvalue \(\lambda\) on \(IH^n\). Sun--Zheng,
Corollary 1.4, p.~237, gives \[
 \mu_{\lambda,e}=\mu_{q^n/\lambda,e},\qquad
 \mu_{q^{n/2},e},\,\mu_{-q^{n/2},e}\in2\mathbf Z
                   \quad\hbox{when }n+e\text{ is even}.
                                                        \tag{CC-XC.P34}
\] In particular, for odd \(n\) each total algebraic multiplicity
\(\sum_e e\mu_{\pm q^{n/2},e}\) is even. The statement explicitly
accommodates nontrivial Jordan blocks. For smooth schemes \(IH^n=H^n\);
for singular proper schemes this theorem remains an intersection-
cohomology theorem, not an assertion of the same purity and parity for
ordinary cohomology.

The properness qualification cannot be discarded from raw odd Betti
parity. For example, \[
 H^1(\mathbf G_{m,\bar k},\mathbf Q_\ell)
      \simeq\mathbf Q_\ell(-1)
                                                        \tag{CC-XC.P35}
\] has dimension one. To see this directly, the Gysin sequence for
\(\mathbf G_m=\mathbf P^1\setminus\{0,\infty\}\) identifies this group
with the kernel of the sum map
\(\mathbf Q_\ell(-1)^2\to\mathbf Q_\ell(-1)\). Thus proper-smooth parity
is the relevant settled question, and it is distinct from statements
about arbitrary open varieties or arbitrary singular proper ordinary
cohomology.

\subsection{CC-XC13: rational spectral targets for proper smooth
schemes}\label{cc-xc13-rational-spectral-targets-for-proper-smooth-schemes}

\textbf{Theorem.} All arithmetic-divisibility and rational-target
conclusions of CC-XC11 hold when ``smooth projective'' is replaced by
``proper smooth''. More explicitly, let \(X/\mathbf F_q\), \(q=p^a\), be
a nonempty proper smooth scheme, pure of dimension \(d\), and write \[
 \det(t-F\mid H^i(X_{\overline{\mathbf F}_q},
                         \mathbf Q_\ell))
             =\prod_f f(t)^{m_f}.
                                                        \tag{CC-XC.P36}
\] For a root \(\alpha_f\) and \(K_f=\mathbf Q(\alpha_f)\), the
finite-field index \(d_f\) in (CC-XC.R15) divides \(m_f\). Hence the
unique rational abelian-sector object \[
 M_i(X)=\bigoplus_f S_f^{\oplus m_f/d_f}
                                                        \tag{CC-XC.P37}
\] has the semisimplification of \(H^i(X)\) as its \(\ell\)-adic
realization for every \(\ell\ne p\). Their sum has a presentation as a
rational summand of a finite direct sum of cohomology motives of abelian
varieties and finite étale schemes, with no Tate twist.

The arithmetic assertion is unconditional. The abelian-sector
interpretation uses the same explicitly scoped all-abelian
classification as CC-XC11; no arbitrary-variety Tate assertion is
introduced.

\subsubsection{Complete proof of the
extension}\label{complete-proof-of-the-extension}

The comparison input needed for the crystalline part also has the
precise proper-smooth extension. Suh, \textbf{Theorem 1.4, p.~296, with
proof on p.~297}, states that the conclusions of his Theorems 1.2 and
1.3 hold for every proper smooth variety over the finite field. Theorem
1.2 is the equality between crystalline and étale Frobenius
characteristic polynomials. Thus, with \(K_0=W(\mathbf F_q)[1/p]\) and
semilinear \(\varphi\) on \(C_i=H^i_{\mathrm{cris}}(X)[1/p]\), \[
 \det(t-\varphi^a\mid C_i)
     =\det(t-F\mid H^i_{\mathrm{et}}(X_{\bar k},
                                      \mathbf Q_\ell)).
                                                        \tag{CC-XC.P38}
\] The same result gives integral coefficients, and proper-smooth purity
gives weight \(i\). The perfectness and Frobenius-invertibility results
in Stacks tags 07MY and 07N5 require proper smoothness, so they apply
directly. Disconnected schemes are handled by the finite component
argument in CC-XC11, or by applying the proper-smooth statements to
their components.

Fix an irreducible \(\mathbf Q_p\)-factor \(f_v\) of one \(f\). The
whole primary block for \(f_v\) is \(\varphi\)-stable exactly as in
(CC-XC.R19), because its polynomial coefficients are fixed by Witt
Frobenius. It has dimension \(m_f[K_{f,v}:\mathbf Q_p]\). The
determinant norm calculation (CC-XC.R20)--(CC-XC.R21), which did not use
projectivity, therefore gives \[
 m_f[K_{f,v}:\mathbf Q_p]\,
          v_p(\alpha_f)/a\in\mathbf Z.
                                                        \tag{CC-XC.P39}
\] All finite-place index denominators divide \(m_f\).

For odd \(i\), use the perfect alternating form (CC-XC.P33) from
Sun--Zheng, Corollary 1.2. Since geometric Frobenius acts on
\(\mathbf Q_\ell(-i)\) by \(q^i\), it is a similitude of multiplier
\(q^i\) for this form. The complete generalized-eigenspace argument
(CC-XC.R22)--(CC-XC.R24) now applies in every odd degree directly: the
generalized eigenspaces at each of \(q^{i/2}\) and \(-q^{i/2}\) have
even dimension. Equivalently, one may use Corollary 1.4, whose more
detailed Jordan-block statement was recorded above.

If \(K_f\) has a real place, \(f\) divides \(t^2-q^i\). Thus for odd
\(i\) the multiplicity \(m_f\) is even, both for a rational real root
and for the real-quadratic factor with its two roots. For even \(i\),
the real-place invariant is zero. Combining this with (CC-XC.P39) proves
\(d_f\mid m_f\).

All remaining arguments of CC-XC11 use only integrality, weights, the
rational abelian-sector classification, and Honda's theorem. They
therefore apply without change: the eigenvalues are \(p\)-units, the
orbit multiplicities descend as in CC-AG08, positive weights are
realized by tensor powers of \(H^1\) after taking a weight-one root, and
degree zero is realized by finite étale schemes. This proves the stated
unique rational target and its untwisted presentation. No semisimplicity
conclusion for \(X\) was used or obtained. \(\square\)

This extension concerns the spectral target. It does not, on its own,
extend the projective intersection-computability and cycle-enumeration
algorithms in CC-XC12 to arbitrary proper nonprojective schemes. The
separate existence of cross-correspondences identifying \(X\) with its
target remains the required arbitrary-variety condition.

\subsection{APA bibliography and precise source
provenance}\label{apa-bibliography-and-precise-source-provenance}

Deligne, P. (1980). La conjecture de Weil: II. \emph{Publications
Mathématiques de l'IHÉS, 52}, 137--252. Corollary 4.1.5 and following
question, printed p.~218.
https://publications.ias.edu/sites/default/files/Number40.pdf

Suh, J. (2012). Symmetry and parity in Frobenius action on cohomology.
\emph{Compositio Mathematica, 148}(1), 295--303.
https://doi.org/10.1112/S0010437X11007056 Theorem 1.4, p.~296, proof
p.~297; Theorem 2.2.2 and Corollary 2.2.3, p.~298; Theorem 3.3.1,
p.~301, proof completed on p.~302; Corollary 3.3.5, p.~302. The
publisher's journal PDF was inspected; the later arXiv upload 2410.01184
is not being treated as a new 2024 result.
https://www.cambridge.org/core/services/aop-cambridge-core/content/view/CE57E10E819362A8F1DB4EC6C7DCBF26/S0010437X11007056a.pdf/symmetry-and-parity-in-frobenius-action-on-cohomology.pdf

Sun, S., \& Zheng, W. (2016). Parity and symmetry in intersection and
ordinary cohomology. \emph{Algebra \& Number Theory, 10}(2), 235--307.
https://doi.org/10.2140/ant.2016.10.235 Theorem 1.1 and Corollary 1.2,
p.~236; Corollaries 1.3--1.4 and Remark 1.5(2), p.~237. The published
primary PDF, including these rendered pages, was inspected directly.
https://msp.org/ant/2016/10-2/ant-v10-n2-p02-s.pdf

The Stacks Project Authors. (2026). \emph{The Stacks Project}. Remark
60.24.13, tag 07MY, and Theorem 60.26.4, tag 07N5.
https://stacks.math.columbia.edu/tag/07MY ;
https://stacks.math.columbia.edu/tag/07N5

Internal proved inputs: CC-XC11, including its complete primary-block
determinant and generalized-eigenspace proofs, and the scoped rational
classification CC-AG08.
